\chapter{Proposed System}

\section{System Overview}

Dobby Studio is strategically conceptualized as a Single Page Application (SPA) delivered as a single HTML file. To entirely circumvent installation friction and hardware dependencies, the system functions seamlessly within any modern web browser environment. Because the application logic executes entirely client-side with minimal server dependency, it maintains exceptional UI performance without requiring the user to own high-end graphic processing hardware. By deploying the environment on the robust Vercel Content Delivery Network (CDN), it achieves global low-latency availability, ensuring both rural weavers and urban designers experience identical performance metrics.

Structurally, the entire application is compartmentalized into three primary, interrelated subsystems operating concurrently. First, the AI Design Engine seamlessly fields natural language text inputs from the user, interacting iteratively with the external edge computation layer to deduce structural definitions. Second, the deterministic Fabric Renderer harnesses mathematical simulation to compile complex sett arrays and boolean structures directly onto an HTML5 Canvas, visualizing everything from microscopic warp interlacements to global drape simulation. Last, the Export System handles the aggregation of the live state matrix into rigidly formatted industrial artifacts, allowing physical fabrication.

\section{Frontend Architecture}

Dobby Studio is implemented as a self-contained Single Page Application (SPA) delivered as a single \texttt{.html} file, requiring no build pipeline, no package manager, and no installation. The application logic executes entirely client-side within a standard web browser, leveraging modern browser APIs to achieve desktop-equivalent performance without requiring high-end hardware.

The visual layout follows a three-panel grid structure: a left-aligned Sidebar contains the manual parameter controls (sett builder, colour picker, weave presets, and physical dimension inputs); the Centre panel hosts the high-fidelity 2D fabric canvas and the 3D drape simulation viewport; the right-aligned Chat Panel serves as the conversational natural language interface. All DOM manipulation and event handling are implemented in \textbf{vanilla JavaScript (ES6+)}, using a centralized mutable state object \texttt{S} that stores the active sett array, weave structure, thread size, EPI, PPI, and repeat count.

State mutations---whether triggered by the sett builder UI, preset selection, or AI inference output---broadcast universally via a master \texttt{updateAll()} dispatcher, enforcing a strict unidirectional data flow without the need for a reactive framework. An ephemeral undo/redo history stack maintains a capped array of recent state snapshots, enabling non-destructive design iteration. Styling is implemented entirely through CSS custom properties, supporting both light and dark themes via a single \texttt{data-theme} attribute toggle on the root element. The application is deployed as a static asset on the Vercel Content Delivery Network (CDN), providing global low-latency availability without a dedicated backend server.

\section{AI Design Pipeline}
The central intelligence capability directing Dobby Studio relies on immediate API interactions utilizing the Groq Cloud infrastructure. The conversational component leverages the highly performant `llama-3.3-70b-versatile` model. When the user dictates a natural language prompt, the backend executes an instantaneous API request mapping the textual input securely against a heavily engineered system prompt explicitly commanding the output of valid text-to-sett JSON configurations.

In order to allow the model to execute intelligent, iterative alterations—such as "make the green stripes twice as wide" or "increase the density of the fabric"—the pipeline utilizes context-aware prompting. The application's current raw state (current sett notation, EPI, PPI parameters, and active weave structure) is silently injected continuously as a system variable payload appending every sequential prompt call, functionally acting as a localized dynamic memory vector without requiring persistent external databases. Should the user lose network connectivity, the system reverts immediately to a localized NLP fallback parser implementing keyword-based mapping algorithms to extract basic colors without incurring API latency.

\section{Fabric Rendering Engine}
Transmuting mathematical variables into accurate visual assets relies on dual rendering parameters. Initially, the engine applies reflective sett expansion computations. Traditional tartan structures do not merely tile; they pivot. The application maps out both the forward array input and functionally reverses and appends the mirrored pivot sequence automatically to simulate industry-standard threading patterns. 

Subsequently, the application loops to generate boolean-based weave matrices. Distinct weave formulas dictate thread visibility. The application provides localized structural algorithms: a plain weave executes a basic alternating 1/1 checkerboard, a 2/2 twill calculates shifts based inversely on index modularity, while 5-end satin constructions map disjointed coordinate interceptions mathematically. By intersecting the colored coordinate points against these dynamically generated boolean matrices, the specific thread color (warp or weft) traversing the topmost visual layer is selected per index coordinate.

To guarantee that the produced graphic accurately simulates a physical, woven artifact rather than a digital spreadsheet, the engine enforces optical color mixing principles inside the cell-by-cell drawing loop. Because physical threads undergo dimensional blending and cast micros-shadows against intersections, the rendering algorithm injects a 70/30 dominant to recessive RGB blend formula locally. To optimize the massive performance overhead caused by iterating across potentially hundreds of thousands of coordinates asynchronously, weave arrays are cached uniquely using a stringified \texttt{\{weave\}\_\{size\}} indexing key, and random trigonometric fiber noise maps are pre-generated during application load to synthesize a physical texture overlay quickly.

\section{Export System}
Ensuring that Dobby Studio outputs manufacturing-grade logic rather than just aesthetic imagery is critical for market success. The Export System generates functional digital artifacts directly from the UI state variables. For generalized creative applications, standard graphical output is achieved organically by forcing the HTML canvas through a `canvas.toBlob()` data-string manipulation, instantly yielding an incredibly high-resolution PNG locally downloaded to the user memory without invoking backend conversions. Similarly, iterative technical metadata variables export natively into portable JSON syntax encapsulating the specific sett geometries, EPI targets, and structural keys.

Conversely, enabling structural mill integration necessitates precise Weave Interchange Format manipulation. The execution logic systematically formats a WIF 1.1 compliant ASCII file layout sequentially mapping standard blocks. The exporter synthesizes crucial headers consisting of `[WIF]` and `[CONTENTS]`, then procedurally maps multidimensional variables generating complex positional coordinates internally across the `[THREADING]`, `[TIEUP]`, `[TREADLING]`, `[WARP]`, and `[WEFT]` arrays. Finally, it dynamically constructs an indexed hexadecimal integer breakdown forming the required `[COLOR TABLE]` required to program individual harness instructions upon older control interfaces.

Finally, the application automatically constructs a styled HTML view natively leveraging the `window.print()` DOM execution generating a comprehensive Spec Sheet PDF. The printout systematically aggregates necessary mill-floor instructions: visually scaled structural thumbnails, extracted color swatches alongside hexadecimal identifiers, formal textile string notations, established physical repeating bounds, and absolute warp ends calculated against an assumed 150-centimeter loom frame utilizing scaling metrics.

\section{Production Specifications Module}
A textile specification sheet is mathematically useless lacking dimensional constraints. The internal production module handles the physical conversion logic, integrating thread proportions alongside precise Ends Per Inch (EPI) and Picks Per Inch (PPI) density indicators provided by the user interface to deduce exact linear measurements directly into centimeters or inches format.

By running calculation blocks resolving dimensional bounds using standard physical formulas ($width = (totalThreads \div EPI) \times 2.54 \text{ cm}$), the environment simulates realistic physical sizing immediately upon the interface parameters altering. To map an aesthetic composition toward a generalized global market fabric width standard (150cm), mathematical adjustments resolve interpolation logic---calculating bounding $warp \text{ } ends = round(150 \div repeatWidthCm \times totalThreads)$---equipping the designer instantaneously with necessary material procurement estimations.

The presented architectural framework outlines the distinct unidirectional data logic implemented within Dobby Studio. Initially, user input---either formatted as text, image, or localized UI manipulations---channels directly into the AI Pipeline or local logic handlers. This orchestrates systemic variable updates against the master state object dictating the sett array structure. Consequently, the updated structural bindings automatically trigger the central mathematical Renderer formulas, instantaneously repainting both the 2D Canvas and refreshing coordinates spanning the 3D physics framework. Ultimately, these static states proceed directly into the Export functions dictating localized user artifact generation without relying on centralized data storage infrastructure. 
